\chapter{曲率}\label{chap:curvature}
\begingroup
\renewcommand{\thefootnote}{}
\footnotetext{\textcopyright\ Springer International Publishing Switzerland 2014.
  L. Godinho and J. Natário, \emph{An Introduction to Riemannian Geometry},
  Universitext. DOI \href{https://doi.org/10.1007/978-3-319-08666-8_4}{10.1007/978-3-319-08666-8\_4}.}
\endgroup

一般 Riemann 流形的局部几何不同于 Euclidean 空间 $\mathbb{R}^n$ 的平直几何. 例如, 2 维球面 $S^2$ (带标准度量) 上测地三角形的内角和总大于 $\pi$. 衡量这种差异的是\textbf{曲率}. 这一概念由 Gauss 在 1827 年的论文\textquotedblleft 曲面的一般研究\textquotedblright 中引入, 又由 Riemann 本人 (于 1854 年) 推广到任意 Riemann 流形. 曲率有许多不同的表现: 测地线偏离的速率, 沿不同向量场的协变导数不交换的程度, 测地三角形内角和与 $\pi$ 的差, 或者将向量沿闭曲线平行移动时该向量旋转的角度.

本章研究曲率的各种刻画和性质. 第~\RGref{sec:curvature}~节引入一般仿射联络的\textbf{曲率算子}, 对 Riemann 流形还引入等价但更具几何性的\textbf{截面曲率}概念. 本节还定义由曲率张量缩并得到的 \textbf{Ricci 曲率张量}和\textbf{标量曲率}. 在广义相对论中, 这些量是表述 Einstein 方程的基础 (第~\RGref{chap:relativity}~章).

第~\RGref{sec:cartan-structure-equations}~节建立 \textbf{Cartan 结构方程}, 这是利用微分形式计算曲率的一种有力方法. 第~\RGref{sec:gauss-bonnet-theorem}~节利用这些方程证明 \textbf{Gauss–Bonnet 定理}, 将紧曲面的曲率与其拓扑联系起来. 这一定理给出一个简单的例子, 说明完备 Riemann 流形的曲率如何制约其拓扑.

第~\RGref{sec:manifolds-of-constant-curvature}~节讨论具有常曲率的完备 Riemann 流形. 这些流形给出弯曲几何的重要例子, 包括 Gauss, Bolyai 和 Lobachevsky 的负曲率非 Euclidean 几何.

最后, 第~\RGref{sec:isometric-immersions}~节研究 Riemann 流形的曲率与其子流形 (带诱导度量) 的曲率之间的关系. 这推广了 Gauss 对曲面的研究, 包括他著名的 \emph{Theorema Egregium}.
